\section{当事人}
\begin{enumerate}
  \item 相对而言是比较复杂的内容。
  \item 三个概念。
        \begin{enumerate}
          \item 诉讼权利能力。以自己的名义起诉应诉的\zd{资格}。自然人/法人/非法人组织。\zd{是适格的基础和前提}。
          \item 诉讼行为能力。\zd{亲自参加}诉讼行使权利履行义务。法人/非法人组织（\zd{依法登记并领取营业执照}）。
                \begin{enumerate}
                  \item 亲自诉讼 \yjt 完全民事行为能力人。
                  \item 法定代理人代为诉讼 \yjt 无/限制民事行为能力人。
                \end{enumerate}
          \item 当事人适格。有作为\zd{本案}的原告被告起诉应诉的\zd{资格}。\zd{结合具体的案件判断}本案的适格当事。一个原则三个例外：
                \begin{enumerate}
                  \item 原则\zd{本案所争议的实体法律关系的双方当事人就是适格当事人}。
                  \item 例外。
                        \begin{enumerate}
                          \item （常见于消极的）确认之诉中对诉讼标的\zd{享有确认利益}的。在实践中不常见。自证清白。
                          \item \zd{享有管理权}的人。失踪人的财产代管人。
                          \item \zd{公益诉讼的机关组织}。
                        \end{enumerate}
                \end{enumerate}

                \begin{kk}{民法典作为实体法所规定的赔偿义务人就是民诉当中的适格被告}
                  常考的。十一个情形。不用去背，结合民法典的规定判断。
                  \begin{enumerate}
                    \item 个人合伙/企业。有没有\zd{依法登记领取营业热照} \yjt 全体合伙人/企业。
                    \item 个体工商户。\zd{先找字号 \yjt 经营者 \yjt 经营者、实际经营者不一致列成共同被告 }
                    \item 无限制民事行为能力人致人损害，和监护人为共同被告。
                    \item 居委、村委、村民小组 \yjt 居委、村委、民小组看有钱没钱。
                    \item 挂靠形式 \yjt 挂靠方、被挂靠方\zd{连带责任}。
                    \item 工作人员 \yjt 工作单位，工作人员绝无可能当事人。
                    \item 企业法人解散。\yi 注销前 \yjt 该法人 \er \zd{未经清算被注销} \yjt 股东、出资人、发起人。
                    \item 劳务派遣被派遣的工作人员致人损害 \yjt 接受派遣的用工单位；派遣单位有过错的相应责任。
                    \item \zd{在一般保证当中，债权人只起诉一般保证人，没有起诉债务人的，法院应当依法追加债务人为共同被告}。
                    \item 死者近亲属当事人地位。
                  \end{enumerate}
                \end{kk}

                \begin{kk}{赔偿义务人}
                  \begin{enumerate}
                    \item 直接规定 A 承担责任，A 为被告。
                    \item A、B \zd{连带责任}。被告\zd{自主选择}，一起告才列为共同被告。
                    \item A 承担责任；B 因过错/管理职责承担相应、补充责任。B 是否作为共同被告看权利人主张。
                  \end{enumerate}
                \end{kk}
        \end{enumerate}
  \item 当事人的变更。
        \begin{enumerate}
          \item 基于法律规定。\yi 自然人死亡 \er 法人合并、分立 \yjt 继承者。
          \item 基于当事人的意思。在讼中实体权利义务发生转移的问题。
                \begin{enumerate}
                  \item 张三起诉李四归还借款十万，官司打到一半，张三就把这十万的债权转让给了王五。
                  \item 继续诉讼，判决对王五产生约束力，追加王五为无独三。可申请替换原告，是否准许法院决定。
                  \item \zd{当事人恒定主义}，没有什么理论上的支撑，技术性操作，为了提高诉讼效率，节约司法资源。
                \end{enumerate}
        \end{enumerate}
\end{enumerate}
